\medskip\noindent\textbf{The modular orbifold and the exact $GT_0^1$ image.}
There is a further exact description of the ambient peripheral
centralizer in the minimal twelve-puncture case.  It also gives
five specified three-cusp tests for the remaining geometric
intersection.  The intersection itself remains unresolved.

Let
\[
 \Delta=\langle a,b\mid a^2=b^3=1\rangle,
 \qquad E=\widehat\Delta=C_2\amalg C_3,
 \qquad d=ab,\quad e=ba.                                \tag{F.7.19}
\]
The isomorphism with the full profinite coproduct follows from
the same universal maps on the two finite factors.
Use the epimorphisms
\[
 \begin{array}{c|cc|c}
     &a&b&\text{kernel}\\ \hline
 q_3:E\to S_3 &(2\ 3)&(1\ 2\ 3)&K\\
 q_5:E\to A_5 &(1\ 2)(3\ 4)&(1\ 3\ 5)&H .
 \end{array}                                             \tag{F.7.20}
\]
The discrete modular orbifold has one order-two point,
one order-three point, and one puncture, whose oriented
loop is $d$.  In the first quotient $d$ has order two,
and in the second it has order five.  The corresponding
torsion-free covers have the following exact data:
\[
 \begin{array}{c|c|c|c|c}
  &\text{degree}&\text{cusp width}&\text{number of cusps}
        &\text{surface type}\\ \hline
  K&6&2&3&(0,3)\\
  H&60&5&12&(0,12).
 \end{array}                                             \tag{F.7.21}
\]
Indeed the quotient maps inject both finite orbifold
stabilizers, so the covers are surfaces.  Their Euler
characteristics are the degree times $-1/6$; the cusp
counts are the indices of $\langle q_i(d)\rangle$ in
the finite quotient.  The resulting genera are zero.
Completion of these finite-index inclusions retains their
full profinite topology.  Thus $K\cong\widehat F_2$,
$H\cong\widehat F_{11}$, and the second extension is
the actual centre-free orbifold extension of the
icosahedral twelve-puncture action.

Here is the precise identification with the specified
action in F.7.  Start with that oriented sphere $S$
and its twelve-vertex puncture set $V$.  Choose a
basepoint $s\in S\setminus V$ outside every nontrivial
stabilizer and use its image in the quotient orbifold.
Choose ordinary positive orbifold loops $u_2,u_3$
around the order-two and order-three points, with
whiskers so that the positive puncture loop is
$u_c=(u_2u_3)^{-1}$.  Set
\[
             a=u_3^{-1}u_2u_3,\qquad b=u_3^{-1};
                   \qquad ab=u_c.
\]
Path lifting gives the monodromy map to the actual
deck group $A_5$.  Its generating pair can be labelled
as the pair $q_5(a),q_5(b)$ in (F.7.20), by the exact
single-orbit calculation below.  This relabels the
deck group and does not change its kernel or the
given subgroup of the sphere mapping class group.
Based loop lifting now identifies the discrete
kernel with $\pi_1(S\setminus V,s)$, and its completion
identifies $H$ with the same specified profinite
surface group.  The chosen lift of the cusp whisker
selects a vertex $v_0\in V$ stabilized by $q_5(d)$.
The actual identification of the twelve labels is
\[
        A_5/\langle q_5(d)\rangle\longrightarrow V,
                   \qquad g\langle q_5(d)\rangle
                                    \longmapsto g v_0.
\]
Under it the cusp $[t d^5t^{-1}]_H$ is exactly the
puncture $q_5(t)v_0$.  In particular the identification
of $H$ carries its pure geometric mapping class
subgroup to the one used in F.7.
Changing the base path changes this identification
by an inner surface-group map; changing the chosen
lift changes it by the specified deck action.
An equivalent choice of the oriented orbifold
presentation gives the corresponding based covering
isomorphism.  These maps conjugate the peripheral
outer group and its geometric subgroup together,
and carry the displayed $A_5$ to the same geometric
action.  They therefore carry the intersection
below to the identical residual centralizer under
that covering identification.

For precision, $x=d^2$ and $y=e^2$ are a full free basis
of $K$.  Schreier rewriting with transversal
$1,a,b,ab,ba,b^{-1}$ gives the following kernel words
$t_g s t_{q_3(gs)}^{-1}$:
\[
 \begin{array}{c|cc}
 t_g&s=a&s=b\\ \hline
 1&1&1\\ a&1&1\\ b&1&1\\
 ab&x&y^{-1}\\ ba&1&y\\ b^{-1}&x^{-1}&1 .
 \end{array}                                             \tag{F.7.22}
\]
Rewriting $a^2$ and $b^3$ at these six representatives
gives only freely trivial words in $x,y$; these twelve
relator checks are included in the certificate.
The resulting Schreier presentation is exactly
$\langle x,y\mid\ \rangle$, giving the stated basis
and then its full profinite basis.  Conjugation by $a$ sends
$(x,y)$ to $(y,x)$, while conjugation by $b$ sends it
to $(y,x^{-1}y^{-1})$.  Writing
$\omega(x)=y$, $\omega(y)=(xy)^{-1}$, the latter
automorphism is exactly $\operatorname{Inn}(y)\omega$.
Therefore the outer
deck action on $K$ is exactly the geometric $S_3$
used in D.35, with these peripheral generators.

\emph{Characteristic kernels and the finite normalization.}
There are six generating order-two/order-three pairs in
$S_3$, forming one orbit under its six inner automorphisms.
There are $120$ such generating pairs in $A_5$, forming
one orbit under conjugation by $S_5$.
For the latter count, fix $(1\ 2)(3\ 4)$: the generating
three-cycles use the fifth letter and one letter from
each of its two pairs, giving eight choices.
Their products with the involution have order five.
The remaining choices generate proper subgroups.
The finite certificate below verifies all $300$ pairs,
as well as the transitive $S_5$ action.
Consequently each epimorphism in (F.7.20) has the unique
kernel of its indicated quotient type.  In particular
both $K$ and $H$ are characteristic in $E$.

We also need the exact quotient-automorphism statement.
The five $V_4$ subgroups of $A_5$ form an intrinsic set.
Its conjugation action is faithful.  An automorphism
acting trivially on these five subgroups acts trivially
on that faithful permutation representation, and hence
on $A_5$.  Thus $\operatorname{Aut}(A_5)$ embeds in $S_5$;
all $120$ conjugations by $S_5$ already give automorphisms,
so this is the full automorphism group.  An odd
conjugation exchanges the two $A_5$ conjugacy classes
of five-cycles.  It follows that an automorphism
preserving the conjugacy class of $q_5(d)$ is inner.
The certificate verifies the five $V_4$ groups,
faithfulness, and both twelve-element five-cycle classes.

\emph{The full peripheral centralizers.}
For $T=K,H$, let $\mathcal P_T^1$ denote the subgroup
of $\operatorname{Out}(T)$ preserving each individual
oriented peripheral element up to conjugacy with
multiplier one.  Put
\[
 \Omega_E^1=\{[\eta]\in\operatorname{Out}(E):
                         \eta(d)\sim_E d\}.
\]
There are canonical continuous isomorphisms, with
the normalizations described below,
\[
 C_{\mathcal P_H^1}(A_5)
       \xrightarrow{\ \sim\ }\Omega_E^1
       \xrightarrow{\ \sim\ }
 C_{\mathcal P_K^1}(S_3)=GT_0^1,
 \qquad GT_0^1=\ker(\lambda:GT_0\to\widehat\Z^\times).
                                                               \tag{F.7.23}
\]
We prove the maps rather than identify the two actions
only by their surface types.

First, for $m=2$ and $T=K$, or for $m=5$ and $T=H$,
\[
                         C_E(d^m)=\overline{\langle d\rangle}
                                      \cong\widehat\Z . 
                                                               \tag{F.7.24}
\]
The element $d^m$ is a primitive cusp of the free group
$T$, so $C_T(d^m)=\overline{\langle d^m\rangle}$.
This follows, for example, by writing $T$ as the
free profinite product of that primitive cyclic factor
and its free complement and using factor malnormality
as in D.35.  The quotient action on the cusp classes
is the actual action on
$Q/\langle q(d)\rangle$, where $Q=S_3$ or $A_5$.
An element centralizing $d^m$ therefore has quotient
in $\langle q(d)\rangle$.  Multiplying it by a suitable
integral power of $d$ puts it in $T$ without changing
centralization.  This proves the equality in (F.7.24).
The procyclic group is a full $\widehat\Z$ because
$d^m$ already generates such a subgroup: a relation
$d^u=1$ would imply $(d^m)^u=1$ and hence $u=0$.

The centralizer $C_E(T)$ is trivial.  Its image in
$Q$ would act trivially in $\operatorname{Out}(T)$,
but that action is faithful on the cusp permutation
set.  Thus this image is trivial, and the remaining
element lies in $Z(T)=1$.
Consequently conjugation embeds $E$ in
$\operatorname{Aut}(T)$ as the full inverse image
of the specified $Q\leq\operatorname{Out}(T)$.

Let $[\phi]\in C_{\mathcal P_T^1}(Q)$.
Conjugation by an actual representative $\phi$ on
$\operatorname{Aut}(T)$ preserves that inverse image
and gives an automorphism $\eta$ of $E$ inducing
the identity on $Q$ and restricting to $\phi$ on $T$.
Changing $\phi$ by an inner automorphism of $T$
changes $\eta$ by the same inner automorphism of $E$.
Also $\eta(d)^m$ is conjugate in $T$ to $d^m$.
After that conjugation, (F.7.24) places $\eta(d)$
in $\overline{\langle d\rangle}$; the equation
$m u=m$ in $\widehat\Z$ forces its exponent to be
one.  Thus $[\eta]\in\Omega_E^1$.

Conversely, choose $[\eta]\in\Omega_E^1$.
Its induced automorphism of $Q$ is inner: for
$S_3$ this follows from its full automorphism group;
for $A_5$ it follows from preservation of the
five-cycle class proved above.
If that quotient automorphism is $\operatorname{Inn}(q_0)$,
choose $u\in E$ with $q(u)=q_0^{-1}$ and replace
$\eta$ by $\eta_0=\operatorname{Inn}(u)\eta$.
The new automorphism induces the identity on $Q$.
Different choices of $u$ differ by an element of $T$,
so the outer class of its restriction to $T$ is
well defined.  Changing the representative of
$[\eta]$ has the same inner effect after normalization;
the centre of $Q$ is trivial, so there is no further
quotient ambiguity.

Write $\eta_0(d)=rdr^{-1}$.  Because the quotient
action is now the identity, $q(r)$ centralizes
$q(d)$.  In both finite groups that centralizer is
exactly $\langle q(d)\rangle$.  Hence $r=t d^j$
for $t\in T$, and $\eta_0(d)=tdt^{-1}$.
For every $s\in E$, the elements $\eta_0(s)$ and
$s$ have the same quotient.  Therefore the restriction
preserves each $T$-cusp class
$[s d^m s^{-1}]_T$ with exponent one.
It belongs to $\mathcal P_T^1$ and centralizes the
outer $Q$ action.

These constructions are inverse homomorphisms.
For injectivity, a normalized automorphism whose
restriction to $T$ is inner can first be adjusted
to fix $T$ pointwise.  Then
$\eta(s)s^{-1}\in C_E(T)=1$ for every $s\in E$,
so it is the identity.  In the other direction,
an inner automorphism of $E$ arising from a
centralizing outer class of $T$ has quotient in
$Z(Q)=1$, hence its conjugator lies in $T$.
Continuity follows from conjugation on the specified
closed subgroup of $\operatorname{Aut}(T)$ and
normalization by the finite quotient.  The maps
are homeomorphisms because all the groups here
are compact Hausdorff.
Finally the equality with $GT_0^1$ is precisely the
Harbater--Schneps centralizer identification used
in D.35, now with the exact basis and outer actions
in (F.7.22)\cite[Propositions 3--4 and Theorem 2]{R-HS}.

Let $\iota:GT_0^1\hookrightarrow\mathcal P_H^1$
denote the inverse of the left-to-right isomorphism
in (F.7.23).  Genus-zero congruence identifies the
completed pure geometric group with its actual
image in $\mathcal P_H^1$.  Thus the remaining group
has the exact description
\[
             C_{\widehat{P\Gamma}_{0,12}}(A_5)
                 =\widehat{P\Gamma}_{0,12}
                                \cap\iota(GT_0^1)
                 \quad\text{inside }\mathcal P_H^1.     \tag{F.7.25}
\]
This proves the relationship with the full $GT_0$
problem; it does not calculate this intersection.

\emph{Explicit substitutions and three-cusp quotient maps.}
Write $\gamma=(1,f)\in GT_0^1$ in its normalized
representative on $K$, so
$x\mapsto x$, $y\mapsto f^{-1}yf$ and $f\in K'$.
Its actual extension to (F.7.19) has the particularly
simple formulas
\[
          \eta_f(a)=f^{-1}a,\qquad
          \eta_f(b)=f^{-1}b,\qquad \eta_f(d)=d .         \tag{F.7.26}
\]
Indeed the first symmetry is $\theta(f)=f^{-1}$.
Direct conjugation on $K$ gives
$\gamma\theta=\operatorname{Inn}(f^{-1})\theta\gamma$,
so $C_E(K)=1$ forces the stated image of $a$.
Since $\gamma(x)=x$, (F.7.24) with $m=2$ forces
$\eta_f(d)=d$; the displayed image of $b=a^{-1}d$
then follows.  These formulas describe the actual
automorphism furnished by (F.7.23), including its
relations and its invertibility.

The quotient automorphism on $A_5$ fixes $q_5(d)$
and is inner.  There is consequently a unique
$j\in\{0,1,2,3,4\}$ for which it is
$\operatorname{Inn}(q_5(d)^j)$; equivalently
\[
             q_5(f)=q_5(e)^j q_5(d)^{-j}.
\]
The pure normalized representative on $H$ is
\[
       \eta_f^H=
          \left.\operatorname{Inn}(d^{-j})\eta_f\right|_H.
                                                               \tag{F.7.27}
\]
For example, with $c_0=d^5$ and $c_1=a d^5a^{-1}=e^5$,
put $h_1=d^{-j}f^{-1}e^j\in H$.  Then
\[
                 c_0\longmapsto c_0,
        \qquad c_1\longmapsto h_1c_1h_1^{-1}.           \tag{F.7.28}
\]
This uses the finite quotient to ensure $h_1\in H$;
it is not a formal substitution of a word outside
the source of a filling homomorphism.

Choose the cusp labels as the twelve left cosets
$A_5/\langle q_5(d)\rangle$, in the breadth-first
order specified by the certificate.  The corresponding
representatives $t_i\in\Delta$ and peripheral words
$c_i=t_i d^5t_i^{-1}$ are recorded there; $t_0=1$
and $t_1=a$.
For a three-element label set $J$ define the actual
continuous filling homomorphism
\[
 q_J:H\longrightarrow H_J
     =H/\langle\!\langle c_i:i\notin J\rangle\!\rangle_H
                         \cong\widehat F_2.            \tag{F.7.29}
\]
The isomorphism type follows from filling precisely
those nine punctures of the specified sphere;
completion preserves the displayed quotient
presentation.  Since (F.7.27) is pure peripheral,
it descends to $H_J$.  Thus there is a continuous
homomorphism
\[
 \mathcal D_J:GT_0^1\longrightarrow\mathcal P_{H_J}^1,
 \qquad\mathcal D_J(\gamma)=
 \bigl[q_J(h)\longmapsto q_J(\eta_f^H(h))\bigr].       \tag{F.7.30}
\]
The brackets on the right specify the outer
automorphism induced by the displayed substitution.
Any change of representative in (F.7.27) changes
it only by an inner automorphism of $H_J$.
Formula (F.7.26), the finite determination of $j$,
and (F.7.29) define each map without assuming that
any two arbitrarily based retained cusp words
form a full free basis.

The pure geometric mapping class group of the
three-punctured sphere is trivial.  The actual
forgetting homomorphism on pure mapping classes,
its action on the filled fundamental group, and
continuity on completions therefore give
\[
 \iota^{-1}\!\left(
       C_{\widehat{P\Gamma}_{0,12}}(A_5)\right)
               \subseteq\bigcap_{|J|=3}\ker\mathcal D_J.
                                                               \tag{F.7.31}
\]
For $g\in A_5$, a lift in $E$ induces the actual
quotient isomorphism $H_J\to H_{gJ}$ by conjugation;
different lifts differ by an inner automorphism.
Because $\iota(\gamma)$ centralizes $A_5$ in
$\operatorname{Out}(H)$, these isomorphisms intertwine
$\mathcal D_J(\gamma)$ and $\mathcal D_{gJ}(\gamma)$.
Their kernels are consequently equal within each
actual $A_5$ orbit of three-element subsets.

The $220$ triples form exactly five orbits, with
the following representatives in the recorded coset
labelling:
\[
 \begin{array}{c|ccccc}
 J&\{0,1,2\}&\{0,1,4\}&\{0,1,8\}&\{0,1,10\}&\{0,4,9\}\\
 |A_5J|&20&60&60&60&20 .
 \end{array}                                             \tag{F.7.32}
\]
The certificate supplies a labelled $A_5$ transporter,
as a word in $a,b$, for every triple.  Hence the
intersection of $220$ kernels in (F.7.31) equals the
intersection of these five specified kernels.
No injectivity or triviality of that intersection
has been proved here.

The standard-library executable
\path{code/part_a/modular_a5_gt0_f7.py} has been run.
It checks the six and $120$ epimorphism pairs and
their single postcomposition orbits, both faithful
cusp actions, both cyclic cusp centralizers, the
five-$V_4$ automorphism argument, the exact twelve
Schreier rewritings in (F.7.22), all twelve peripheral
words, and all $220$ subset transporters in (F.7.32).
These finite calculations certify the finite inputs
and the displayed words.  The centre-free extension,
normalization, completion, and filling maps are
proved above.
